\documentclass[12pt, a4paper]{article}
\usepackage[utf8]{inputenc}
\usepackage{amsmath, amssymb, amsthm}
\usepackage{geometry}

\geometry{margin=1in}

% Theorem Environments
\newtheorem{theorem}{Theorem}
\newtheorem{lemma}{Lemma}
\theoremstyle{definition}
\newtheorem{definition}{Definition}

% Custom Commands
\newcommand{\Sem}{\mathcal{S}} % Semi-algebra
\newcommand{\Alg}{\mathcal{A}(\mathcal{S})} % Generated Algebra
\newcommand{\Union}{\bigcup}
\newcommand{\DisjointUnion}{\bigsqcup}

\title{Proof of the Extension Theorem from a Semi-Algebra to an Algebra}
\author{}
\date{}

\begin{document}

\maketitle

\section*{Problem Statement}
Let $\Omega$ be a set and let $\Sem \subseteq \mathcal{P}(\Omega)$ be a semi-algebra. Let $\mu: \Sem \to [0, \infty]$ be a $\sigma$-additive set function such that $\mu(\emptyset) = 0$.

We wish to prove:
\begin{enumerate}
    \item \textbf{Existence:} There exists a set function $\nu: \Alg \to [0, \infty]$ which extends $\mu$ (i.e., $\nu(E) = \mu(E)$ for all $E \in \Sem$) and is $\sigma$-additive on $\Alg$.
    \item \textbf{Uniqueness:} If $\nu_1$ and $\nu_2$ are two such extensions, then $\nu_1(A) = \nu_2(A)$ for all $A \in \Alg$.
\end{enumerate}

\section{Preliminaries}
Recall the structure of the algebra generated by a semi-algebra $\Sem$.
\begin{lemma}
    If $\Sem$ is a semi-algebra, then the algebra $\Alg$ generated by $\Sem$ consists of all finite disjoint unions of sets in $\Sem$. That is, for every $A \in \Alg$, there exist pairwise disjoint sets $S_1, \dots, S_n \in \Sem$ such that:
    \[
    A = \DisjointUnion_{i=1}^n S_i
    \]
\end{lemma}

\section{Part 1: Existence}

\subsection{Step 1: Definition of $\nu$}
For any $A \in \Alg$, let $A = \DisjointUnion_{i=1}^n S_i$ be a representation of $A$ with $S_i \in \Sem$. We define $\nu$ as:
\begin{equation} \label{def:nu}
    \nu(A) := \sum_{i=1}^n \mu(S_i)
\end{equation}
Since $\mu(\emptyset)=0$, adding empty sets to the representation does not change the sum.

\subsection{Step 2: Well-Definedness}
We must show that $\nu(A)$ is independent of the representation of $A$. Suppose $A$ has two representations:
\[
A = \DisjointUnion_{i=1}^n S_i = \DisjointUnion_{j=1}^m T_j, \quad \text{where } S_i, T_j \in \Sem.
\]
Consider the intersection $S_i \cap T_j$. Since $\Sem$ is a semi-algebra, $S_i \cap T_j \in \Sem$.
Fix $i$. We can write $S_i$ as:
\[
S_i = S_i \cap A = S_i \cap \left( \DisjointUnion_{j=1}^m T_j \right) = \DisjointUnion_{j=1}^m (S_i \cap T_j)
\]
Since $\mu$ is $\sigma$-additive (and thus finitely additive) on $\Sem$, and the sets $(S_i \cap T_j)_{j=1}^m$ are disjoint elements of $\Sem$, we have:
\[
\mu(S_i) = \sum_{j=1}^m \mu(S_i \cap T_j)
\]
Summing over all $i=1, \dots, n$:
\[
\sum_{i=1}^n \mu(S_i) = \sum_{i=1}^n \sum_{j=1}^m \mu(S_i \cap T_j)
\]
By symmetry, performing the same operation on $T_j = \DisjointUnion_{i=1}^n (T_j \cap S_i)$ yields:
\[
\sum_{j=1}^m \mu(T_j) = \sum_{j=1}^m \sum_{i=1}^n \mu(T_j \cap S_i)
\]
The double sums are identical. Thus, $\sum \mu(S_i) = \sum \mu(T_j)$, proving $\nu$ is well-defined.
\textit{Note: Finite additivity of $\nu$ on $\Alg$ follows immediately from this definition.}

\subsection{Step 3: $\sigma$-Additivity of $\nu$}
Let $A \in \Alg$ and let $\{A_k\}_{k=1}^\infty$ be a sequence of disjoint sets in $\Alg$ such that $A = \DisjointUnion_{k=1}^\infty A_k$. We must show $\nu(A) = \sum_{k=1}^\infty \nu(A_k)$.

Since $A \in \Alg$, we can decompose it into elementary sets $S_1, \dots, S_n \in \Sem$:
\[
A = \DisjointUnion_{i=1}^n S_i
\]
For each $k$, since $A_k \in \Alg$, usually $A_k$ is a union of sets in $\Sem$. However, we can simply observe that $(S_i \cap A_k) \in \Alg$. Thus, $(S_i \cap A_k)$ can be written as a finite disjoint union of sets in $\Sem$:
\[
S_i \cap A_k = \DisjointUnion_{j=1}^{m_{ik}} R_{ikj}, \quad \text{where } R_{ikj} \in \Sem
\]
Now, observe the decomposition of the base set $S_i$:
\[
S_i = S_i \cap A = S_i \cap \left( \DisjointUnion_{k=1}^\infty A_k \right) = \DisjointUnion_{k=1}^\infty (S_i \cap A_k) = \DisjointUnion_{k=1}^\infty \DisjointUnion_{j=1}^{m_{ik}} R_{ikj}
\]
The set $S_i$ is now written as a \textbf{countable disjoint union} of sets $R_{ikj} \in \Sem$. Since $\mu$ is $\sigma$-additive on $\Sem$:
\[
\mu(S_i) = \sum_{k=1}^\infty \sum_{j=1}^{m_{ik}} \mu(R_{ikj})
\]
By the definition of $\nu$, the inner sum corresponds to $\nu(S_i \cap A_k)$:
\[
\sum_{j=1}^{m_{ik}} \mu(R_{ikj}) = \nu\left( \DisjointUnion_{j=1}^{m_{ik}} R_{ikj} \right) = \nu(S_i \cap A_k)
\]
Substituting this back:
\[
\mu(S_i) = \sum_{k=1}^\infty \nu(S_i \cap A_k)
\]
Finally, we sum over $i$ to get $\nu(A)$:
\begin{align*}
    \nu(A) &= \sum_{i=1}^n \mu(S_i) \\
           &= \sum_{i=1}^n \sum_{k=1}^\infty \nu(S_i \cap A_k) \\
           &= \sum_{k=1}^\infty \sum_{i=1}^n \nu(S_i \cap A_k) \quad (\text{Valid to swap finite and infinite sums})
\end{align*}
Observe the inner sum $\sum_{i=1}^n \nu(S_i \cap A_k)$. Since $\{S_i \cap A_k\}_{i=1}^n$ is a finite disjoint decomposition of $A_k$, and $\nu$ is finitely additive (from Step 2), this sum equals $\nu(A_k)$.
\[
\nu(A) = \sum_{k=1}^\infty \nu(A_k)
\]
Thus, $\nu$ is $\sigma$-additive.

\section{Part 2: Uniqueness}

Let $\nu_1$ and $\nu_2$ be two $\sigma$-additive measures on $\Alg$ that extend $\mu$.
Let $A \in \Alg$. By the structure of the algebra, there exist disjoint sets $S_1, \dots, S_n \in \Sem$ such that:
\[
A = \DisjointUnion_{i=1}^n S_i
\]
Since $\nu_1$ is $\sigma$-additive on $\Alg$, it is finitely additive. Thus:
\[
\nu_1(A) = \sum_{i=1}^n \nu_1(S_i)
\]
By the hypothesis, $\nu_1$ extends $\mu$, so $\nu_1(S_i) = \mu(S_i)$ for all $S_i \in \Sem$. Therefore:
\[
\nu_1(A) = \sum_{i=1}^n \mu(S_i)
\]
Applying the same logic to $\nu_2$:
\[
\nu_2(A) = \sum_{i=1}^n \nu_2(S_i) = \sum_{i=1}^n \mu(S_i)
\]
Comparing the two results:
\[
\nu_1(A) = \sum_{i=1}^n \mu(S_i) = \nu_2(A)
\]
Thus, $\nu_1(A) = \nu_2(A)$ for all $A \in \Alg$. The extension is unique.

\end{document}